\setchapterpreamble[o]{%
  \chapterdictum[Dummy Person]{%
    Dummy quote from Dummy Person.%
  }%
}
\chapter{Dummy Chapter}\label{main:dummy}

This chapter briefly showcases the main structural elements, environments, and
helper commands provided by the template. Much of this functionality comes from
existing packages, either through direct use or through the particular
configuration adopted here. The same is true for more general aspects such as
layout, headings, and floats, where the template relies primarily on
KOMA-Script, in particular its \texttt{scrbook} class. For further details, see
the files in the \texttt{preamble} directory, which are heavily commented and
contain links to the relevant packages and their documentation.

\section{Textual features}\label{main:dummy:text}

This section demonstrates ordinary prose features. For example, this sentence
contains a regular footnote\footnote{This is a footnote.} and a margin
note.\marnote{This is a margin note.} Margin notes take two compilation passes
to properly typeset at the right location. This sentence contains two
consecutive footnotes, separated by a comma, inserted using
\verb|\multiplefootnoteseparator|.%
\footnote{The first of two consecutive footnotes.}\multiplefootnoteseparator%
\footnote{The second of two consecutive footnotes.}

When drafting (that is, when \verb|\ifdraft| evaluates to true), todo notes with
different urgency levels are available as well:
\begin{itemize}
\item Very\todoA{High-priority todo note.} urgent.

\item Less urgent.\todoB{Medium-priority todo note.}

\item Optional,\todoC{Low-priority todo note.} perhaps.
\end{itemize}

All todo notes are automatically added to the list of todo notes at the
beginning of the document.

\subsection{Subsection heading}\label{main:dummy:text:sectioning} Subsection
headings are the deepest heading level that is still numbered and included in
the table of contents.

\subsubsection{Subsubsection heading} Subsubsection headings are no longer
numbered and are not included in the table of contents.

\paragraph{Paragraph heading} Paragraph headings are run in and end with a
period.

\subparagraph{Subparagraph heading} Subparagraph headings are similar to
paragraph headings, but are also indented.

\section{Math}\label{main:dummy:math}

This section showcases some of the template's built-in mathematical
typesetting, including theorem-like environments, displayed equations, and a
few convenience macros.

\begin{definition}[Sampling]\label{main:dummy:def:sample}
  For a set \(S\), the expression \(x \sample S\) denotes that \(x\) is sampled from
  \(S\) according to some distribution.
\end{definition}

\begin{remark}[Sampling macro]\label{main:dummy:rmk:sample}
  In \cref{main:dummy:def:sample}, the sampling operator is typeset using the
  \verb|\sample| macro provided by the template.
\end{remark}

\begin{axiom}[Toy axiom]\label{main:dummy:ax:toy}
  For every positive integer \(n\), there exists a string in \(\{0,1\}^n\).
\end{axiom}

\begin{lemma}[Toy lemma]\label{main:dummy:lem:toy}
  Let \(n \ge 1\). If \(x \sample \{0,1\}^n\), then \(x \in \{0,1\}^n\).
\end{lemma}
\begin{proof}
  Immediate from the definition of sampling from \(\{0,1\}^n\).
\end{proof}

\begin{proposition}[Toy proposition]\label{main:dummy:prop:toy}
  Let \(x \sample \{0,1\}^n\). Then \(x\) has length \(n\).
\end{proposition}

\begin{theorem}[Toy theorem]\label{main:dummy:thm:toy}
  Every positive integer admits a binary string of the same length.
\end{theorem}
\begin{proof}[Proof sketch]
  Choose the string guaranteed by \cref{main:dummy:ax:toy}.
\end{proof}

\begin{corollary}\label{main:dummy:cor:toy}
  For every positive integer \(n\), the set \(\{0,1\}^n\) is non-empty.
\end{corollary}

\begin{claim}\label{main:dummy:clm:counter}
  \Cref{main:dummy:ax:toy,main:dummy:lem:toy,main:dummy:prop:toy,main:dummy:thm:toy,main:dummy:cor:toy}
  share the same theorem counter.
\end{claim}

Below is an example of a displayed equation:
\begin{equation}\label{main:dummy:eq:identity}
  (a+b)^2 = a^2 + 2ab + b^2
\end{equation}

As \cref{main:dummy:eq:identity} illustrates, equation tags are placed in the
margin, thereby keeping them visually separate from the mathematical content and
freeing up space for the typesetting of the equation itself.

\section{Figures and Tables}\label{main:dummy:figtab}

The template aims to provide sensible default formatting for figures, tables,
and the corresponding lists of figures and tables. Caption labels
are typeset in bold to distinguish them from the caption text; single-line
captions are centered, while multi-line captions are set hanging from the second
line onward. By default, figure captions are configured for placement below the
figure, while table captions are configured for placement above the table, in
line with common typographic convention. The actual caption placement nevertheless
remains under user control: placing \verb|\caption| before the content puts the
caption above it, while placing it after the content puts it below.
See \cref{main:dummy:fig:ex,main:dummy:tab:ex} for examples.

\begin{figure}
  \centering
  \colorbox{T-Q-B0}{\rule[4em]{4em}{0pt}}
  \caption{Example figure with single-line caption.}
  \label{main:dummy:fig:ex}
\end{figure}

\begin{table}
  \centering
  \caption{Example table with an artificially long caption to ensure that it
  spans multiple lines and thus demonstrates the hanging indentation used for
  longer captions.}
  \label{main:dummy:tab:ex}
  \begin{tabular}{lcc}
    \toprule
    Foo & Bar & Baz  \\
    \midrule
    Foo cell & Bar cell & Baz cell \\
    Foo cell & Bar cell & Baz cell \\
    Foo cell & Bar cell & Baz cell \\
    \bottomrule
  \end{tabular}
\end{table}

\section{Algorithms and Listings}\label{main:dummy:alglst}

Beyond figures and tables, the template also provides floating environments for
algorithms and listings, together with matching content lists. These are styled
consistently with the other floating environments (and their lists). As
for tables, the captions for algorithms and listings are formatted
for placement above the content.

For algorithms, the template uses a floating environment that does not add a
frame around the algorithm contents by itself. Because of that, the contents
need to provide their own framing. To make this easy, the template provides the
\verb|algorithmicframed| environment. It works like \verb|algorithmic|, but
automatically adds a top and bottom rule of the correct width. In addition, the
template changes some of the default keywords and formatting used by the
algorithm environments make the result a bit less cluttered/verbose. See
\cref{main:dummy:alg:ex} for an example.

\begin{algorithm}
  \centering
  \caption{Example algorithm.}
  \label{main:dummy:alg:ex}
  \begin{algorithmicframed}[1]
    \Procedure{Random}{\(n\)}
    \LComment{Line comment}
    \For{\(i = 0, \ldots, 10\)}
      \State \(x \sample \{0,1\}^n\) \Commentvs{Comment (default space)}
    \EndFor
    \State \Return \(x\) \Commentvs[3em]{Comment (custom space)}
    \EndProcedure
  \end{algorithmicframed}
\end{algorithm}

For users of \verb|cryptocode|, the template also adjusts and enables that
package’s typesetting to better match the other pseudocode environments used
here (\verb|algorithmic| and \verb|algorithmicframed|). This mainly
affects the font size, line numbers, line spacing, and comment formatting. See
the example in \cref{main:dummy:cc:ex}, and compare the typesetting of its code
and line numbers with that of \cref{main:dummy:alg:ex}. The goal is not to make
the two styles identical, but to make their main typographic elements
sufficiently coherent that they work well together.

\begin{figure}
  \begin{pchstack}[center,boxed]
    \procedure[linenumbering]{$\mathsf{Random}(n)$}{%
      \pclinecomment{Line comment}\\
      \pcfor i = 0, \ldots, 10\pccolon\\
      \t x \sample \{0,1\}^n \pccomment{Comment (default space)}\\
      \pcreturn x \pccomment[3em]{Comment (custom space)}
    }
  \end{pchstack}
  \caption{Example (\texttt{cryptocode}) procedure.}
  \label{main:dummy:cc:ex}
\end{figure}

Finally, the template provides a default setup for the \verb|listings| package.
This is intended for source-code listings written using \verb|lstlisting|
environments. Although \verb|listings| also provides its own support for
floating, captions, labels, and content lists, you may prefer to use the
template's \verb|listing| environment instead, so that these elements remain
consistent with the other floating environments. See \cref{main:dummy:lst:ex}
for an example.
\begin{listing}
  \caption{Example listing.}
  \label{main:dummy:lst:ex}
\begin{lstlisting}[language={[3]Python}]
def random(n: int):
    for i in range(0, 11):
        x = os.urandom(n)
    return x
\end{lstlisting}
\end{listing}

\section{Colors}\label{main:dummy:colors}

By default, the template loads the \verb|colorblind| package. Besides providing
several colorblind-friendly color schemes, this package also remaps a number of
the standard built-in color names to a colorblind-friendly default scheme. At
the same time, several built-in color names for which no suitable default
colorblind-friendly replacement is provided are intentionally remapped to black.
\Cref{main:dummy:tab:colors} shows these remapped default color names as used by
the template. If you want to use \verb|colorblind| without this default
remapping behavior, pass it the \verb|keep-defaults| option.

\begin{table}
  \newcommand{\swatch}[1]{%
    \fcolorbox{black}{#1}{\rule{0pt}{1.2ex}\rule{2.5em}{0pt}}%
  }
  \centering
  \caption{Remapped built-in color names provided by \texttt{colorblind}.}
  \label{main:dummy:tab:colors}
  \begin{tabular}{llc}
    \toprule
    Name & Swatch & RGB \\
    \midrule
    \texttt{blue}      & \swatch{blue}      & (68, 119, 170) \\
    \texttt{cyan}      & \swatch{cyan}      & (102, 204, 238) \\
    \texttt{green}     & \swatch{green}     & (34, 136, 51) \\
    \texttt{yellow}    & \swatch{yellow}    & (204, 187, 68) \\
    \texttt{red}       & \swatch{red}       & (238, 102, 119) \\
    \texttt{violet}    & \swatch{violet}    & (170, 51, 119) \\
    \texttt{gray}      & \swatch{gray}      & (187, 187, 187) \\
    \midrule
    \texttt{orange}    & \swatch{orange}    & (0, 0, 0) \\
    \texttt{brown}     & \swatch{brown}     & (0, 0, 0) \\
    \texttt{darkgray}  & \swatch{darkgray}  & (0, 0, 0) \\
    \texttt{lightgray} & \swatch{lightgray} & (0, 0, 0) \\
    \texttt{lime}      & \swatch{lime}      & (0, 0, 0) \\
    \texttt{magenta}   & \swatch{magenta}   & (0, 0, 0) \\
    \texttt{olive}     & \swatch{olive}     & (0, 0, 0) \\
    \texttt{pink}      & \swatch{pink}      & (0, 0, 0) \\
    \texttt{purple}    & \swatch{purple}    & (0, 0, 0) \\
    \texttt{teal}      & \swatch{teal}      & (0, 0, 0) \\
    \bottomrule
  \end{tabular}
\end{table}

\section{Helpers}\label{main:dummy:helpers}

The template also provides a few small helpers for inline and mathematical
typesetting. For example, to typeset a text subscript and superscript at the
same time, it provides \verb|\textsubsuperscript|:
H\textsubsuperscript{2}{+}O.

For outlined text, the template provides \verb|\tbox| (solid), \verb|\tdbox|
(dashed), and \verb|\tdtbox| (dotted):

\begin{itemize}
\item \tbox{Text, solid outline} \tbox[2pt]{Text, more spacious solid outline}
\item \tdbox{Text, dashed outline} \tdbox[2pt]{Text, more spacious dashed outline}
\item \tdtbox{Text, dotted outline} \tdtbox[2pt]{Text, more spacious dotted outline}
\end{itemize}

Similarly, for outlined mathematics, the template provides
\verb|\tboxed| (solid), \verb|\tdboxed| (dashed), and \verb|\tdtboxed|
(dotted). These commands are context-aware, meaning that they scale with the
surrounding math style, for example in subscripts and superscripts:

\begin{itemize}
\item
  \(\tboxed{x+y}\) \(\tboxed[2pt]{x+y}\)
  \(A^{\tboxed{ij}}\) \(A^{\tboxed[2pt]{ij}}\)
  \(B_{\tboxed{k}}\) \(B_{\tboxed[2pt]{k}}\)
\item
  \(\tdboxed{x+y}\) \(\tdboxed[2pt]{x+y}\)
  \(A^{\tdboxed{ij}}\) \(A^{\tdboxed[2pt]{ij}}\)
  \(B_{\tdboxed{k}}\) \(B_{\tdboxed[2pt]{k}}\)
\item
  \(\tdtboxed{x+y}\) \(\tdtboxed[2pt]{x+y}\)
  \(A^{\tdtboxed{ij}}\) \(A^{\tdtboxed[2pt]{ij}}\)
  $B_{\tdtboxed{k}}$ $B_{\tdtboxed[2pt]{k}}$
\end{itemize}
